\section{Proof Appendices}
\label{sec:appendices}

The proofs below rely on the structural separation results of Section~\ref{sec:separation}~\cite{F02,F03,F04,F08,K01,K02,S01,S02}.

\subsection{Proof of Theorem~\ref{thm:T11} (Energy-optimal pairwise input)}
Let $Q = \Delta H^* R^{-1} \Delta H$ be the Hermitian, nonnegative definite matrix arising from the linearized output sensitivity $\Delta H$ and measurement noise covariance $R$. The average output power for an input $u(t)$ with power spectral density $\Phi_u(\omega)$ is $\frac{1}{2\pi}\int_{-\infty}^{\infty} \Phi_u(\omega) d\omega$, while the output signal-to-noise ratio in the frequency domain is $\Phi_u(\omega) \cdot \lambda_{\max}(Q(\omega))/\text{tr}(R)$. Maximizing the output SNR under a total power constraint $P$ leads to the Rayleigh-Ritz quotient $\max_{\|u\|_2^2\leq P} u^* Q u / u^* u = \lambda_{\max}(Q)$, whose maximizer is the eigenvector corresponding to $\lambda_{\max}(Q)$. In the frequency domain this yields the optimal input spectrum $\Phi_u^*(\omega) = P \cdot v_{\max}(\omega) v_{\max}^*(\omega)$, where $v_{\max}(\omega)$ is the principal eigenvector of $Q(\omega)$. The time-domain optimal input is thus the principal eigenfunction of the operator $Q$, and the optimal amplitude is $\sqrt{P\lambda_{\max}(Q)}$. The boxed formulas for $D_{\max}$ and $u^*$ follow directly from Parseval's theorem and the definition of $Q$.

\subsection{Proof of Theorem~\ref{thm:T12} (Finite multisineness sufficiency)}
Let the objective be $J(\phi) = \int_{-\pi}^{\pi} \Phi_u(e^{i\omega}) \cdot w(\omega) d\omega$, where $w(\omega)\ge 0$ is a trigonometric polynomial of degree $P$ representing the frequency-weighted discriminant (e.g., $\lambda_{\max}(Q(\omega))$). By the Fejer-Riesz theorem, any nonnegative trigonometric polynomial of degree $P$ can be written as $|A(e^{i\omega})|^2$ for some polynomial $A(z)$ of degree at most $P$. Hence $J(\phi) = \frac{1}{2\pi}\int_{-\pi}^{\pi} |A(e^{i\omega})|^2 \Phi_u(e^{i\omega}) d\omega$. By the Krein-Milman theorem, the maximum of a linear functional over the convex set of power spectra with total power $P$ is attained at an extreme point, which are Dirac combs. Since $A(z)$ has at most $P$ zeros on the unit circle, the optimal spectrum is supported at most at $P+1$ points (the $P$ zeros plus possibly one extra point where the polynomial does not vanish). Thus a multisine with at most $P+1$ distinct frequencies suffices.

\subsection{Proof of Theorem~\ref{thm:T13} (Optimal pairwise sampling schedule)}
Consider the set of candidate sampling frequencies $\{\omega_1,\dots,\omega_N\}$ with associated weights $w_i = \lambda_{\max}(Q(\omega_i))$. The goal is to select $P$ frequencies to maximize $\sum_{i\in\mathcal{S}} w_i$ subject to a minimum separation constraint $\Delta$. This is a classic weighted interval scheduling problem. Define $f(k)$ as the optimal total weight for the first $k$ frequencies sorted in ascending order. The recurrence $f(k)=\max\{f(k-1),\, w_k + f(p(k))\}$ where $p(k)$ is the last index $<k$ with $\omega_k-\omega_{p(k)}\ge\Delta$ yields the optimal solution via dynamic programming in $O(N\log N)$ time. The additive exchange argument shows that any optimal schedule can be transformed into the greedy top-$n$ schedule when the separation constraint is inactive, proving the top-$n$ rule for the unconstrained case.

\subsection{Proof of Theorem~\ref{thm:T16} (Existence of a safe and informative perturbation)}
Let $\mathcal{U}_\rho = \{u\in L_\infty[0,T] : \|u\|_\infty \le \rho\}$ be the set of admissible inputs bounded by the safety margin $\rho>0$. The input-to-state map $\Phi: u\mapsto x(t)$ is continuous in the $L_\infty$ norm for the fractional-order system under consideration (by the Lipschitz continuity of the Caputo operator and the Lipschitz vector field). The information gain $\mathcal{I}(u)$ is a continuous functional of the state trajectory (being an integral of quadratic forms of the state). Hence the composition $\mathcal{I}\circ\Phi$ is continuous on the compact, convex set $\mathcal{U}_\rho$. By the extreme value theorem, a maximum exists. Moreover, if $\rho$ is chosen smaller than the distance from the nominal equilibrium to the unsafe boundary, the image of $\Phi$ stays inside the safe region, guaranteeing safety of the maximizing input.

\subsection{Proof of Theorem~\ref{thm:T17} (Fractional inward-pointing rectangle invariant)}
Consider the rectangle $\mathcal{R}=[x_L,x_U]\times[y_L,y_U]$ and the Caputo fractional system ${}^C\!D_t^\alpha z = f(z)+Bu$. Let $V(x,y)=\max\{x_L-x,\,x-x_U,\,y_L-y,\,y-y_U\}$ be the distance to the boundary of $\mathcal{R}$. Using the comparison principle for Caputo derivatives (see e.g., Li \& Yao, 2018) and the fact that $f$ satisfies the quasimonotone condition on the faces of $\mathcal{R}$ (due to the strong Allee inequalities), one obtains ${}^C\!D_t^\alpha V \le -\eta V$ for some $\eta>0$ whenever $V>0$. Comparison with the scalar fractional differential equation ${}^C\!D_t^\alpha w = -\eta w$ yields $V(t)\le E_\alpha(-\eta t^\alpha)V(0)$, where $E_\alpha$ is the Mittag-Leffler function, which decays to zero. Hence trajectories starting inside $\mathcal{R}$ remain inside for all $t\ge0$.

\subsection{Proof of Theorem~\ref{thm:T18} (Nonlinear Bayesian posterior update)}
Given the linear-Gaussian approximation of the observation model around the current posterior mean $\bar\theta_n$, the likelihood is $p(y_{n+1}|M,\theta_{M})\approx \mathcal{N}(H_M\theta_{M},R)$. The joint posterior factorizes as $p(M,\theta_M|y_{1:n+1})\propto p(y_{n+1}|M,\theta_{M})p(\theta_M|M,y_{1:n})p(M|y_{1:n})$. Integrating out $\theta_M$ yields a Gaussian posterior for each model, and the model posterior updates via Bayes' rule with marginal likelihood $\mathcal{N}(y_{n+1};H_M\bar\theta_{M|n}, H_M P_{M|n} H_M^\top+R)$. The gain $Q=\Delta H^* R^{-1}\Delta H$ appears when computing the KL divergence between prior and posterior model distributions, confirming the reduction to the linear-Gaussian case of Section~\ref{sec:design}.

\subsection{Proof of Theorem~\ref{thm:T19} (Bayesian sequential information bound)}
Let $p_n(M)=p(M|y_{1:n})$ and $p_{n+1}(M)=p(M|y_{1:n+1})$. The conditional mutual information satisfies $I(M;y_{n+1}|y_{1:n},d_{n+1}) = \mathbb{E}_{p(d_{n+1})}[D_{\mathrm{KL}}(p_{n+1}(\cdot)\|p_n(\cdot))]$. By the concavity of entropy, $H(M|y_{1:n+1})\ge H(M|y_{1:n}) - I(M;y_{n+1}|y_{1:n},d_{n+1})$, which rearranges to $I(M;y_{n+1}|y_{1:n},d_{n+1}) \le H(M|y_{1:n})$. Equality holds iff the posterior collapses to a point mass, i.e., a single experiment suffices to identify the model. The chain rule for mutual information yields $I(M;y_{1:N}) = \sum_{k=1}^N I(M;y_k|y_{1:k-1})$, so the expected number of experiments needed to reduce the entropy to zero is at least $\lceil H(M)/\max I\rceil$. The greedy policy that maximizes the one-step information gain therefore minimizes this bound.